\corrigesChapitre{13}{Circuits : schématiser et protéger}
\par\noindent\begin{minipage}{\linewidth}
\begin{corrige}{1}{Une boucle}
\begin{questions}
\item Non, dans le circuit simple étudié.
\item Il fournit de l’énergie au circuit.
\end{questions}
\pcFigureEx{figures/cor-circuit-ouvert.pdf}{L’ouverture interrompt l’unique boucle ; on ne dessine pas un courant permanent à travers l’espace ouvert.}
\end{corrige}
\end{minipage}
\par\noindent\begin{minipage}{\linewidth}
\begin{corrige}{2}{Le sens conventionnel}
\begin{questions}
\item De la borne positive du générateur.
\item Vers la borne négative, par le trajet extérieur fermé.
\end{questions}
\pcFigureEx{figures/cor-sens-courant.pdf}{Les flèches indiquent le sens conventionnel dans le circuit extérieur : de + vers −.}
\end{corrige}
\end{minipage}
\par\noindent\begin{minipage}{\linewidth}
\begin{corrige}{3}{En série}
\begin{questions}
\item Elles sont en série.
\item Elle interrompt le courant dans les deux lampes.
\end{questions}
\pcFigureEx{figures/cor-serie-ouverte.pdf}{Une ouverture dans la seule boucle interrompt le courant dans les deux lampes.}
\end{corrige}
\end{minipage}
\par\noindent\begin{minipage}{\linewidth}
\begin{corrige}{4}{En dérivation}
\begin{questions}
\item Une association en dérivation.
\item Oui ; c’est précisément le cas du montage étudié.
\end{questions}
\pcFigureEx{figures/cor-branches.pdf}{Entre les mêmes deux nœuds : une branche de générateur et deux branches de lampes.}
\end{corrige}
\end{minipage}
\par\noindent\begin{minipage}{\linewidth}
\begin{corrige}{5}{La protection}
\begin{questions}
\item Un court-circuit, notamment aux bornes du générateur.
\item Non : il doit interrompre le circuit lorsqu’un dépassement approprié se produit, selon ses caractéristiques.
\end{questions}
\end{corrige}
\end{minipage}
\par\noindent\begin{minipage}{\linewidth}
\begin{corrige}{6}{Compter les branches}
\begin{questions}
\item Trois branches : générateur, lampe 1, lampe 2.
\item Deux branches contiennent chacune une lampe.
\end{questions}
\pcFigureEx{figures/cor-branches.pdf}{Entre les mêmes deux nœuds : une branche de générateur et deux branches de lampes.}
\end{corrige}
\end{minipage}
\par\noindent\begin{minipage}{\linewidth}
\begin{corrige}{7}{Placer l’interrupteur}
\begin{questions}
\item Dans la branche de la lampe 1, en série avec elle.
\item Sur un trajet commun entre le générateur et le point de séparation, ou dans sa branche, selon le schéma.
\end{questions}
\par\smallskip\noindent\textbf{Pour comprendre.}\enspace Suivre un trajet fermé pour chaque lampe. Placer l’interrupteur dans le trajet propre à L1 ne coupe pas celui de L2 ; le placer dans le fil commun interrompt les deux trajets.\par
\pcFigureEx{figures/cor-commandes.pdf}{Un interrupteur dans la branche L1 commande seulement L1 ; un interrupteur dans le trajet commun peut commander les deux.}
\end{corrige}
\end{minipage}
\par\noindent\begin{minipage}{\linewidth}
\begin{corrige}{8}{Une lampe dévissée}
\begin{questions}
\item L2 peut continuer de fonctionner.
\item Le générateur reste adapté, l’autre branche intacte et l’ouverture ne coupe pas une connexion commune.
\end{questions}
\pcFigureEx{figures/cor-branche-ouverte.pdf}{La branche L1 est ouverte ; L2 garde un trajet fermé vers la source, supposée adaptée et stable.}
\end{corrige}
\end{minipage}
\par\noindent\begin{minipage}{\linewidth}
\begin{corrige}{9}{Suivre les connexions}
\begin{questions}
\item Comparer les nœuds et les dipôles reliés à chaque paire de nœuds.
\item Non : la topologie des connexions compte, pas l’apparence du tracé.
\end{questions}
\par\smallskip\noindent\textbf{Pour comprendre.}\enspace La forme du fil peut être déformée sans changer le circuit tant que ses connexions restent les mêmes. Repérer les bornes de chaque dipôle puis suivre la boucle dans les deux dessins.\par
\pcFigureEx{figures/cor-circuits-equivalents.pdf}{La position sur la feuille change, mais la même succession de dipôles relie les bornes de la source.}
\end{corrige}
\end{minipage}
\par\noindent\begin{minipage}{\linewidth}
\begin{corrige}{10}{Un fil parasite}
\begin{questions}
\item Il court-circuite la lampe.
\item Le courant peut la contourner ; la répartition des tensions et l’intensité peuvent changer. L’effet dépend du reste du circuit.
\end{questions}
\pcFigureEx{figures/cor-lampe-court-circuitee.pdf}{Le fil ajouté relie directement les bornes de L1 ; ne pas tester ce défaut sur une pile.}
\end{corrige}
\end{minipage}
\par\noindent\begin{minipage}{\linewidth}
\begin{corrige}{11}{Un cahier des charges}
\begin{questions}
\item Les placer en dérivation sur des branches séparées, avec une alimentation adaptée.
\item Non : un défaut commun au générateur, aux fils partagés ou aux protections peut couper tout le circuit.
\end{questions}
\pcFigureEx{figures/cor-branche-ouverte.pdf}{La branche L1 est ouverte ; L2 garde un trajet fermé vers la source, supposée adaptée et stable.}
\end{corrige}
\end{minipage}
\par\noindent\begin{minipage}{\linewidth}
\begin{corrige}{12}{Un défaut à ne pas tester}
\begin{questions}
\item Le courant et l’échauffement peuvent devenir dangereux et endommager la pile.
\item Analyser un schéma ou un dispositif scolaire protégé conduit par le professeur.
\end{questions}
\end{corrige}
\end{minipage}
\par\noindent\begin{minipage}{\linewidth}
\begin{corrige}{13}{Une protection et sa limite}
\begin{questions}
\item Il faut suivre les consignes de sécurité et faire examiner le défaut par une personne compétente ; l’élève ne manipule pas l’installation.
\item Les protections ont des fonctions et des conditions de déclenchement différentes ; elles ne suppriment pas toute possibilité de contact dangereux.
\end{questions}
\end{corrige}
\end{minipage}
\par\noindent\begin{minipage}{\linewidth}
\begin{corrige}{14}{Une autonomie trompeuse}
\begin{questions}
\item Non : le trajet partagé a été interrompu.
\item L’indépendance concerne l’ouverture d’une branche réceptrice, avec générateur et connexions communes conservés.
\end{questions}
\pcFigureEx{figures/cor-fil-commun.pdf}{L’indépendance des branches suppose que le générateur et les connexions communes fonctionnent.}
\end{corrige}
\end{minipage}
\par\noindent\begin{minipage}{\linewidth}
\begin{corrige}{15}{Relier explication et schéma}
\begin{questions}
\item Sa description correspond à une série, pas au schéma en dérivation.
\item Décrire la séparation du courant aux nœuds et les deux trajets récepteurs, ou redessiner une véritable association en série.
\end{questions}
\pcFigureEx{figures/cor-courants-paralleles.pdf}{Le courant principal se partage au nœud ; les lampes ne sont pas traversées successivement par le même courant.}
\end{corrige}
\end{minipage}
\par\noindent\begin{minipage}{\linewidth}
\phantomsection\label{cor-defi-13}\pcSolutionDefi{13}{L’éclairage du refuge}\begin{questions}
\item Relier chaque lampe à une branche entre les bornes du générateur ; placer l’interrupteur en série avec L2 seulement.
\item Seule la branche L2 est ouverte ; L1 reste alimentée dans le modèle.
\item Un fil commun coupé, le générateur ou une protection commune peuvent interrompre tout le montage.
\end{questions}
\pcFigureEx{figures/cor-refuge.pdf}{L1 est directement alimentée ; seul l’interrupteur de la branche L2 commande L2.}
\end{minipage}
\par\noindent\begin{minipage}{\linewidth}
\phantomsection\label{cor-activite-13}\pcSousTitre{Pistes pour l’investigation}
Placer les deux lampes dans deux branches dérivées et l’interrupteur en série avec une seule. L’ouverture coupe uniquement cette branche si source et connexions communes restent intactes. Le schéma doit montrer les deux nœuds et un trajet fermé pour l’autre lampe. Le matériel et les tensions doivent être adaptés ; la vérification précède l’alimentation.\par\end{minipage}
